\documentclass[11pt]{article}

\usepackage[margin=1.2in]{geometry}
\usepackage{amsmath,amssymb}
\usepackage{microtype}
\usepackage{parskip}
\usepackage{listings}
\usepackage{xcolor}
\usepackage{hyperref}
\usepackage{booktabs}
\usepackage{array}
\usepackage{enumitem}
\usepackage{mdframed}

\definecolor{codegray}{RGB}{245, 245, 245}
\definecolor{codecomment}{RGB}{100, 116, 139}
\definecolor{codestring}{RGB}{22, 101, 52}
\definecolor{codenumber}{RGB}{100, 116, 139}
\definecolor{bordercolor}{RGB}{203, 213, 225}
\definecolor{accentblue}{RGB}{30, 64, 175}
\definecolor{warningyellow}{RGB}{254, 243, 199}
\definecolor{warningborder}{RGB}{217, 119, 6}
\definecolor{notegreen}{RGB}{220, 252, 231}
\definecolor{noteborder}{RGB}{22, 101, 52}

\lstdefinelanguage{Rust}{
  keywords={pub, struct, enum, fn, let, mut, impl, use, self, Self,
            match, if, else, for, while, return, true, false,
            Vec, Option, Result, HashMap, HashSet, Some, None, Ok, Err,
            type, where, trait, as, in, ref, mod, super, Box, Arc,
            String, bool, u32, u64, f64, usize},
  keywordstyle=\color{accentblue}\bfseries,
  comment=[l]{//},
  commentstyle=\color{codecomment}\itshape,
  stringstyle=\color{codestring},
  string=[b]",
  sensitive=true,
  morecomment=[s]{/*}{*/},
}

\lstset{
  language=Rust,
  backgroundcolor=\color{codegray},
  basicstyle=\ttfamily\footnotesize,
  breaklines=true,
  keepspaces=true,
  showstringspaces=false,
  frame=single,
  framerule=0.5pt,
  rulecolor=\color{bordercolor},
  xleftmargin=12pt,
  xrightmargin=12pt,
  aboveskip=8pt,
  belowskip=8pt,
  numbers=left,
  numberstyle=\tiny\color{codenumber},
  numbersep=8pt,
  escapeinside={(*@}{@*)},
}

\newcommand{\entity}[1]{\texttt{\textbf{#1}}}
\newcommand{\field}[1]{\texttt{#1}}

\newenvironment{prosbox}{%
  \begin{mdframed}[backgroundcolor=notegreen,
    linecolor=noteborder, linewidth=1pt,
    innertopmargin=8pt, innerbottommargin=8pt,
    innerleftmargin=10pt, innerrightmargin=10pt]
  \small
}{%
  \end{mdframed}
}

\newenvironment{consbox}{%
  \begin{mdframed}[backgroundcolor=warningyellow,
    linecolor=warningborder, linewidth=1pt,
    innertopmargin=8pt, innerbottommargin=8pt,
    innerleftmargin=10pt, innerrightmargin=10pt]
  \small
}{%
  \end{mdframed}
}

\title{Prescription Modeling in the FEN Schema:\\
Design Options and Medication Interactions}
\date{\today}

\begin{document}
\maketitle

\section{Problems with the Current Prescription Payload}
\label{sec:current-problems}

The current \entity{Prescription} payload in the FEN schema is:

\begin{lstlisting}
Prescription {
    medication:    CodedValue,
    dosage:        Dosage,
    frequency:     Frequency,     // currently Vec<Duration>
    period:        Option<TimeInterval>,
    relevant_note: Option<FactId>,
}
\end{lstlisting}

Three issues with this representation motivate the design options
explored in the rest of this document.

\subsection{Issue 1: \texttt{Frequency} as \texttt{Vec<Duration>}}
\label{sec:freq}

As currently typed, \field{Frequency} is a \texttt{Vec<Duration>}~---
a sequence of time gaps between doses. This representation has
significant limitations.

For a regular schedule, all entries in the vector are identical:
three gaps of eight hours to express TID, two gaps of twelve hours for
BID. The repetition carries no additional information. More
importantly, the representation cannot express:

\begin{itemize}[noitemsep]
  \item \field{PRN} (as needed)~--- no regular interval exists
  \item Frequency conditioned on meals, symptoms, or lab values
    (``take with food,'' ``hold if systolic below 90'')
  \item Titration schedules (``0.25mg weekly for four weeks,
    then 0.5mg weekly'')
  \item Standard named frequencies used universally in clinical
    practice (QD, BID, TID, QID, QHS)
\end{itemize}

It is also misaligned with standard medical coding systems that
already handle this well: FHIR's \texttt{Timing} type, the RxNorm
term structure, and NCI Thesaurus frequency codes all provide richer,
standardized representations.

A typed enum addresses these limitations:

\begin{lstlisting}
pub enum Frequency {
    // Standard named frequencies
    QD,         // once daily
    BID,        // twice daily
    TID,        // three times daily
    QID,        // four times daily
    QHS,        // at bedtime
    QOD,        // every other day
    Weekly,
    BiWeekly,
    Monthly,
    // Interval-based for non-standard schedules
    EveryNHours { n: u32 },
    EveryNDays  { n: u32 },
    EveryNWeeks { n: u32 },
    // Conditional -- no fixed interval
    PRN { indication: Option<String> },
    // Complex or titration schedules
    Complex { description: String, coded: Option<CodedValue> },
}
\end{lstlisting}

This change is independent of the lifecycle design options discussed
in Section~\ref{sec:lifecycle} and applies regardless of which option
is adopted.

\subsection{Issue 2: Temporal Representation --- \texttt{occurred\_at} and \texttt{period}}
\label{sec:temporal}

\field{Fact.occurred\_at} is a \field{TemporalAnchor}, which can be
either a \field{Point(Timestamp)} or a
\field{Period\{start, end\}}. The \entity{Prescription} payload also
carries \field{period: Option<TimeInterval>}. This creates two
overlapping temporal fields with no clear contract for what each means
or which takes precedence.

For discrete events~--- lab results, imaging results, encounters~---
\field{occurred\_at} is unambiguous. A lab result happened at a point;
an encounter spanned a period. One field covers both cases cleanly.

For prescriptions the situation is more complex:

\paragraph{Under the current single-snapshot model.}
The prescription is written at a point in time (\field{occurred\_at}
as \field{Point}), but it remains active for an unknown duration.
A clinician writes a prescription without knowing when it will be
stopped or modified. The actual duration of the prescription course is
therefore an open period~--- start known, end unknown.
\field{TemporalAnchor::Period} cannot express this: it requires both
\field{start} and \field{end}. The \field{period} field in the payload
was presumably added to capture the intended course length (``for 10
days''), but this is distinct from the actual duration and creates two
temporal fields that mean different things.

\paragraph{Under the event-per-fact model (Option~A or Option~B).}
This ambiguity dissolves. Each prescription event is a discrete
occurrence: the \field{Start} event happened at a point in time (when
the prescription was written); the \field{Stop} event happened at a
point in time (when it was discontinued). \field{occurred\_at} is
always \field{Point} for prescription events. The duration of the
prescription course is not a property of any single fact~--- it is
derived by the \field{FoldMedicationState} evaluator as the span
between the \field{Start} and \field{Stop} events for a given
\field{medication\_order\_id}. The \field{period} payload field is no
longer needed and can be removed.

One remaining edge case: a prescription written for a defined finite
course (``amoxicillin 500mg TID for 10 days'') has a known
\emph{intended} end date at the time of writing, which differs from
the actual stop date if the patient stops early or the prescriber
extends the course. Under the event-per-fact model this is captured as
an \field{intended\_duration: Option<Duration>} on the \field{Start}
event, not as a period timestamp:

\begin{lstlisting}
pub enum PrescriptionEvent {
    Start {
        intended_duration: Option<Duration>,  // None for indefinite
    },
    Modify  { changed: Vec<PrescriptionField> },
    Stop    { reason: Option<String> },
    Hold    { reason: Option<String>,
              expected_resume: Option<ApproximateDate> },
    Resume,
    Renew,
}
\end{lstlisting}

\field{intended\_duration} is enrichable by the LLM path from note
text and is distinct from the actual course duration.

\paragraph{Under Option~D (separate top-level entity).}
\entity{MedicationOrder} carries \field{started\_at} and
\field{stopped\_at} directly as first-class fields. The temporal
ambiguity in the \entity{Fact} model does not arise because the
medication order is no longer a \entity{Fact}.

\subsection{Issue 3: The Lifecycle Gap}
\label{sec:lifecycle-gap}

The current payload captures a single snapshot. It does not record
what changed, why, or how a prescription relates to the one it
replaced. A dosage reduction, a route change, a temporary hold, and a
renewal are all clinically distinct events~--- but the current model
treats them identically: a new \entity{Prescription} fact with
different field values, with no structured record of the transition.

The options in Section~\ref{sec:lifecycle} address this gap directly.

\section{Prescription Lifecycle Options}
\label{sec:lifecycle}

\subsection{Option A: Event-per-Fact Within the Existing Enum}
\label{sec:option-a}

Keep \entity{Prescription} as a \entity{FactPayload} variant, but add
an explicit event type field that records what kind of prescription
event this fact represents. Each lifecycle change is a new \entity{Fact}
with a stable \field{medication\_order\_id} that groups all events for
the same medication course.

\begin{lstlisting}
pub enum PrescriptionEvent {
    Start { intended_duration: Option<Duration> },
    Modify {
        changed: Vec<PrescriptionField>,
    },
    Stop {
        reason: Option<String>,
    },
    Hold {
        reason:           Option<String>,
        expected_resume:  Option<ApproximateDate>,
    },
    Resume,
    Renew,
}

pub enum PrescriptionField {
    Dosage, Frequency, Route,
}

Prescription {
    medication:           CodedValue,
    medication_order_id:  OrderGroupId,
    event:                PrescriptionEvent,
    dosage:               Dosage,
    frequency:            Frequency,
    route:                Option<String>,
    rationale:            Option<String>,
    relevant_note:        Option<FactId>,
}
\end{lstlisting}

The \field{medication\_order\_id} is stable across all prescription
events for a single medication course, providing a direct grouping key
without requiring chain traversal. The \field{period} payload field is
removed: course duration is derived from the \field{Start} and
\field{Stop} events. The \field{intended\_duration} on \field{Start}
handles finite prescribed courses.

\paragraph{Example.}

\begin{lstlisting}
// Initial prescription (2022)
Fact {
    id: FactId("f-met-001"),
    occurred_at: TemporalAnchor::Point(Timestamp("2022-03-01")),
    payload: Prescription {
        medication:          CodedValue { system: RxNorm, code: "860975",
                                         display: "Metformin 1000mg" },
        medication_order_id: OrderGroupId("mo-metformin-chen"),
        event: PrescriptionEvent::Start {
            intended_duration: None,  // indefinite
        },
        dosage:       Dosage { amount: 1000.0, unit: "mg" },
        frequency:    Frequency::BID,
        rationale:    None,
        relevant_note: None,
    },
    ..
}

// Dose reduction (2026-01-15)
Fact {
    id: FactId("f-met-002"),
    occurred_at: TemporalAnchor::Point(Timestamp("2026-01-15")),
    payload: Prescription {
        medication:          CodedValue { system: RxNorm, code: "860975",
                                         display: "Metformin 500mg" },
        medication_order_id: OrderGroupId("mo-metformin-chen"),
        event: PrescriptionEvent::Modify {
            changed: vec![PrescriptionField::Dosage,
                          PrescriptionField::Frequency],
        },
        dosage:       Dosage { amount: 500.0, unit: "mg" },
        frequency:    Frequency::QD,
        rationale:    Some("Dose reduced due to declining renal \
                            function, eGFR approximately 38"),
        relevant_note: Some(FactId("f-note-patel-20260115")),
    },
    ..
}
\end{lstlisting}

\begin{prosbox}
\textbf{Pros.}
\begin{itemize}[noitemsep]
  \item Stays within the \entity{Fact} enum --- each event is genuinely
    a discrete ground truth event from Epic
  \item \field{medication\_order\_id} provides a direct grouping key;
    no chain traversal needed
  \item \field{FoldMedicationState} evaluator processes the ordered
    list of events for a course to compute current active medications
  \item Deterministic path maps cleanly: Epic's order modification
    events have explicit modification type codes that map directly to
    \field{PrescriptionEvent} variants
  \item \field{occurred\_at} is always \field{Point}~--- the temporal
    ambiguity of the current model is resolved
  \item \field{period} payload field is no longer needed
  \item Task~0 note linking and Task~1 enrichment apply naturally
\end{itemize}
\end{prosbox}

\begin{consbox}
\textbf{Cons.}
\begin{itemize}[noitemsep]
  \item Querying the current dose requires selecting the most recent
    prescription event for the course, not reading a single record
  \item The \field{changed} field in \field{Modify} is informational
    but not enforced: the full dosage/frequency must still be repeated
    even when only one changed
  \item \field{medication\_order\_id} is a new identifier type not
    in the original schema; it must be assigned and kept stable across
    encounters at deterministic uplift time
\end{itemize}
\end{consbox}

\subsection{Option B: Separate PrescriptionOrder and PrescriptionEvent Variants}
\label{sec:option-b}

Split into two distinct \entity{FactPayload} variants.
\entity{PrescriptionOrder} captures the initial order as written.
\entity{PrescriptionEvent} captures each discrete lifecycle change,
referencing the order it modifies. Only the fields that actually
changed are populated on the event.

\begin{lstlisting}
// The initial order -- written once
PrescriptionOrder {
    medication:          CodedValue,
    dosage:              Dosage,
    frequency:           Frequency,
    route:               Option<String>,
    intended_duration:   Option<Duration>,
    relevant_note:       Option<FactId>,
}

// Each lifecycle change -- carries only what changed
PrescriptionEvent {
    order_fact_id:  FactId,
    event_type:     PrescriptionEventType,
    new_dosage:     Option<Dosage>,
    new_frequency:  Option<Frequency>,
    new_route:      Option<String>,
    reason:         Option<String>,
    relevant_note:  Option<FactId>,
}

pub enum PrescriptionEventType {
    DosageChange,
    FrequencyChange,
    RouteChange,
    Stop,
    Hold { expected_resume: Option<ApproximateDate> },
    Renew,
    Resume,
}
\end{lstlisting}

\begin{prosbox}
\textbf{Pros.}
\begin{itemize}[noitemsep]
  \item Clean separation: the order is the standing record; events are
    mutations of it
  \item \entity{PrescriptionEvent} carries only fields that changed~---
    no repetition of unchanged dosage or frequency
  \item Querying the event history for a medication is a direct lookup
    by \field{order\_fact\_id}
  \item \field{occurred\_at} is always \field{Point} on both fact types
\end{itemize}
\end{prosbox}

\begin{consbox}
\textbf{Cons.}
\begin{itemize}[noitemsep]
  \item \field{FoldMedicationState} must join across two fact types to
    compute current state, rather than processing a uniform event stream
  \item The \entity{PrescriptionOrder} fact has an unusual lifecycle:
    it is never \field{Superseded} in the traditional sense; its
    evolution is expressed entirely through events
  \item More complex deterministic uplift: Epic's data must be mapped
    to a two-type structure rather than a single event stream
\end{itemize}
\end{consbox}

\subsection{Option C: Reuse the Superseded Chain}
\label{sec:option-c}

Do not add an event type. Each prescription change creates a new
\entity{Prescription} fact with the updated fields, and the prior
prescription receives \field{FactStatus::Superseded\{reason:
ClinicalRefinement\}}.

\begin{lstlisting}
Prescription {
    medication:    CodedValue,
    dosage:        Dosage,
    frequency:     Frequency,
    route:         Option<String>,
    stop_reason:   Option<String>,
    relevant_note: Option<FactId>,
}
\end{lstlisting}

\begin{prosbox}
\textbf{Pros.}
\begin{itemize}[noitemsep]
  \item Simplest representation --- no new fields or types
  \item Reuses existing \field{FactStatus} machinery
  \item Current active medication query: find \field{Active}
    prescription facts for this subject
\end{itemize}
\end{prosbox}

\begin{consbox}
\textbf{Cons.}
\begin{itemize}[noitemsep]
  \item \textbf{Semantic conflation.} \field{Superseded} was designed
    for AI-enrichment and clinical refinement of the \emph{same event}.
    A dosage change is a new clinical action, not a refinement of
    a prior fact. Using \field{Superseded} for prescription lifecycle
    changes corrupts its semantics and makes it impossible to
    distinguish AI enrichment supersessions from clinical modification
    supersessions.
  \item No event type information: a dosage change, a route change,
    and a renewal are indistinguishable without comparing field values
    across versions
  \item The reason for a change has no natural home
  \item The temporal ambiguity of the current model is not resolved:
    \field{occurred\_at} would still need to represent an open period
    for the active prescription
\end{itemize}
\end{consbox}

\subsection{Option D: Prescription as a Separate Top-Level Entity}
\label{sec:option-d}

Promote \entity{MedicationOrder} out of the \entity{Fact} enum
entirely, making it a first-class entity analogous to
\entity{ProblemEpisode}. Prescription lifecycle events become a
timeline on the \entity{MedicationOrder}.

\begin{lstlisting}
pub struct MedicationOrder {
    pub id:           MedicationOrderId,
    pub subject_id:   SubjectId,
    pub medication:   CodedValue,
    pub status:       MedicationOrderStatus,
    pub started_at:   Timestamp,
    pub stopped_at:   Option<Timestamp>,  // None while active
    pub authored_by:  Author,
    pub provenance:   Provenance,
}

pub enum MedicationOrderStatus {
    Active,
    OnHold,
    Stopped  { reason: Option<String> },
    Completed,
}

pub struct MedicationOrderEvent {
    pub id:            EventId,
    pub order_id:      MedicationOrderId,
    pub occurred_at:   Timestamp,
    pub event_type:    MedicationOrderEventType,
    pub new_dosage:    Option<Dosage>,
    pub new_frequency: Option<Frequency>,
    pub reason:        Option<String>,
    pub provenance:    Provenance,
}

pub enum MedicationOrderEventType {
    Start, DosageChange, FrequencyChange, RouteChange,
    Hold, Resume, Stop, Renew,
}
\end{lstlisting}

The temporal ambiguity does not arise under this model: the medication
order carries \field{started\_at} and \field{stopped\_at} as
first-class fields. An active medication order has
\field{stopped\_at: None}, which is a natural representation of an
open-ended course with no sentinel value required.

\begin{prosbox}
\textbf{Pros.}
\begin{itemize}[noitemsep]
  \item Cleanest semantic model: a medication order is more like an
    episode (a persistent evolving object) than a lab result (a point
    observation)
  \item Status and history are first-class properties; current
    medication query is \field{MedicationOrder} where
    \field{status == Active}
  \item \field{stopped\_at: Option<Timestamp>} naturally represents
    an open-ended course without sentinel values
  \item Enables \entity{MedicationRelation} as a typed edge between
    medication orders (see Section~\ref{sec:interactions})
  \item Aligns with FHIR's \texttt{MedicationRequest} model
\end{itemize}
\end{prosbox}

\begin{consbox}
\textbf{Cons.}
\begin{itemize}[noitemsep]
  \item Breaks the \entity{Fact} enum's guarantee. Prescription acts
    (written, modified, stopped) are discrete ground truth events from
    Epic; promoting \entity{MedicationOrder} to a separate entity
    makes it a derived summary rather than a raw observation.
  \item Introduces a fourth top-level layer: Facts, Episodes,
    Narratives, and MedicationOrders. Every part of the system~---
    query engine, security model, read models~--- must learn about a
    new entity type.
  \item \entity{EpisodeMembership} currently links \entity{Fact}
    records to episodes. Medication orders would need their own
    membership mechanism, or the membership model must be generalized.
  \item The \field{Authored} and \field{Temporal} traits apply less
    uniformly across the schema.
\end{itemize}
\end{consbox}

\section{Other ``Dynamic'' Fact Types}
\label{sec:dynamic}

The prescription discussion raises a broader question: which other
\entity{FactPayload} variants share this ``dynamic'' character?

\paragraph{Coverage.}
A patient enrolls in a plan, changes plans, has a gap, and re-enrolls.
The current \entity{Coverage} payload captures a single enrollment
period. The same event-type modeling applies: each enrollment, plan
change, or termination is a separate fact with an event type and a
\field{coverage\_group\_id} to link events for the same coverage
relationship.

\begin{lstlisting}
Coverage {
    payer:              String,
    plan:               String,
    coverage_group_id:  CoverageGroupId,
    event:              CoverageEvent,   // Enroll, PlanChange,
                                         // Terminate, Renew
    member_id:          String,
    coverage_type:      CoverageType,
}
\end{lstlisting}

Note that \field{period} is removed here too: the coverage period
is derived from the \field{Enroll} and \field{Terminate} events
for a given \field{coverage\_group\_id}, mirroring the prescription
approach.

\paragraph{PreAuthRequest and PreAuthDecision.}
These already model part of a lifecycle (request then decision), but
the lifecycle continues: denial is appealed, approval is modified,
authorization expires and requires renewal. The \entity{Appeal} variant
handles some of this, but there is no \field{authorization\_group\_id}
tying the whole lifecycle together. Adding one would allow ``show me
the full authorization history for this procedure request'' as a direct
lookup.

\paragraph{Diagnosis.}
Certainty changes (working $\to$ confirmed $\to$ ruled out) are
enrichments on the same diagnostic assertion, handled by the LLM
enrichment path populating \field{certainty}. A new diagnostic
assertion superseding a prior one (reclassifying a condition,
correcting a misdiagnosis) is genuinely a new event, handled correctly
by \field{FactStatus::Superseded\{reason: ClinicalRefinement\}}.
Diagnosis does not need event-type modeling.

\textcolor{blue}{We are removing certainty modeling / tracking from Diagnosis.}

\section{Medication Interactions}
\label{sec:interactions}

\subsection{What Is Currently Missing}

Medication interactions are not modeled anywhere in the current schema.
They appear only as:

\begin{itemize}[noitemsep]
  \item \field{DecisionPoint.rationale}~--- free text, not queryable
  \item \field{Alternative}~--- captures drugs not chosen, not the
    relationship between drugs that were chosen together
  \item \entity{EpisodeRelation}~--- captures problem-to-problem
    relationships but not medication-to-medication relationships
\end{itemize}

Cohort queries like ``show me all patients on a counter-medication
for a drug they are currently taking'' are not answerable from
structured data.

\subsection{Taxonomy of Medication Relationships}

``Interaction'' covers several clinically distinct concepts:

\begin{enumerate}[noitemsep]
  \item \textbf{Contraindication / conflict.} Drug A and Drug B should
    not be taken together; a conscious override was made and documented.
  \item \textbf{Counter-medication.} Drug B is prescribed to manage a
    side effect of Drug A. Proton pump inhibitor alongside NSAIDs;
    semaglutide for renal protection alongside metformin.
  \item \textbf{Synergistic / joint regimen.} Drugs A and B are
    intentionally prescribed together because they work better in
    combination.
  \item \textbf{Dose-modifying interaction.} Drug A changes the
    metabolism of Drug B, requiring a dose adjustment (e.g.\ warfarin
    and many co-medications).
  \item \textbf{Causal chain.} Drug A causes a condition that requires
    Drug B. NSAIDs $\to$ GI damage $\to$ omeprazole.
  \item \textbf{Sequential / superseding.} Drug B replaces Drug A
    because A was ineffective, caused intolerable side effects, or the
    clinical picture changed.
\end{enumerate}

\subsection{Option I: MedicationRelation as a Standalone Struct}
\label{sec:interaction-option-i}

A typed edge between medication orders, paralleling how
\entity{EpisodeRelation} works between \entity{ProblemEpisode} records.
Most natural under Option~D where the edge links
\entity{MedicationOrder} records; possible but more awkward under
Options~A or~B where it would link \entity{FactId} values representing
prescription events rather than medication courses.

\begin{lstlisting}
pub struct MedicationRelation {
    pub id:               RelationId,
    pub source_order_id:  MedicationOrderId,
    pub target_order_id:  MedicationOrderId,
    pub relation_type:    MedicationRelationType,
    pub asserted_by:      Author,
    pub asserted_at:      Timestamp,
    pub provenance:       Provenance,
}

pub enum MedicationRelationType {
    CounterMedication {
        managed_effect: Option<CodedValue>,
    },
    CoPrescribed {
        rationale: Option<String>,
    },
    Supersedes {
        reason: MedicationSupersessionReason,
    },
    DoseModifying {
        direction: DoseAdjustmentDirection,
    },
    ContraindicationOverride {
        clinical_justification: Option<String>,
    },
    CausedNeedFor,
}

pub enum MedicationSupersessionReason {
    Ineffective,
    SideEffect { effect: Option<CodedValue> },
    ClinicalProgressionChange,
    PatientPreference,
    Other,
}

pub enum DoseAdjustmentDirection { Increase, Decrease }
\end{lstlisting}

\paragraph{Example --- Margaret's medications.}

\begin{lstlisting}
MedicationRelation {
    source_order_id: MedicationOrderId("mo-semaglutide"),
    target_order_id: MedicationOrderId("mo-metformin"),
    relation_type: MedicationRelationType::CoPrescribed {
        rationale: Some("Added for glycemic control and renal \
                         protection alongside reduced metformin"),
    },
    asserted_by: Author { author_type: AuthorType::AiAssisted, .. },
    ..
}
\end{lstlisting}

\begin{prosbox}
\textbf{Pros.}
\begin{itemize}[noitemsep]
  \item Fully queryable as structured data
  \item Authored, provenance-bearing, and contestable
  \item Can be produced by the LLM path from note text or by a
    drug-drug interaction knowledge base at prescription time
\end{itemize}
\end{prosbox}

\begin{consbox}
\textbf{Cons.}
\begin{itemize}[noitemsep]
  \item Most natural under Option~D; linking \entity{FactId} values
    under Options~A or~B is semantically awkward (edge between events
    rather than courses)
  \item Interaction data not stored in structured form in Epic; must
    be asserted by LLM or knowledge base
  \item Adds a third parallel edge type alongside \entity{EpisodeRelation}
    and \entity{NarrativeRelation}
\end{itemize}
\end{consbox}

\subsection{Option II: Encode Interactions Through FactRole on EpisodeMembership}
\label{sec:interaction-option-ii}

Extend \field{FactRole} with interaction-aware roles. The relationship
between two medications is implicit in the roles they play within the
same episode.

\begin{lstlisting}
pub enum FactRole {
    TriggeringSymptom,
    DiagnosticTest,
    Treatment,
    OutcomeMeasure,
    Monitoring,
    Complication,
    Referral,
    Administrative,
    InsuranceAction,
    CounterMedication,
    CoPrescribed,
    Other,
}
\end{lstlisting}

\paragraph{Example.} Naproxen causes GI symptoms; omeprazole manages them.

\begin{lstlisting}
EpisodeMembership { fact_id: f-naproxen,
                    episode_id: ep-arthritis,
                    role: FactRole::Treatment, .. }
EpisodeMembership { fact_id: f-naproxen,
                    episode_id: ep-gi-symptoms,
                    role: FactRole::TriggeringSymptom, .. }
EpisodeMembership { fact_id: f-omeprazole,
                    episode_id: ep-gi-symptoms,
                    role: FactRole::CounterMedication, .. }
\end{lstlisting}

The relationship between naproxen and omeprazole is implicit via the
shared episode, not a direct typed edge.

\begin{prosbox}
\textbf{Pros.}
\begin{itemize}[noitemsep]
  \item No new entity types or relation structs
  \item Small, backward-compatible \field{FactRole} additions
  \item Consistent with the existing episode membership mechanism
\end{itemize}
\end{prosbox}

\begin{consbox}
\textbf{Cons.}
\begin{itemize}[noitemsep]
  \item Relationship between medications is implicit, not queryable
    directly (``what medication is omeprazole countering?'' requires
    joining across episode memberships)
  \item Conflates the role within an episode with the relationship
    between medications
  \item Cannot express dose-modifying interactions or contraindication
    overrides
  \item Cross-episode interactions are particularly hard to represent
\end{itemize}
\end{consbox}

\subsection{Option III: Structured MedicationDecisionContext on DecisionPoint}
\label{sec:interaction-option-iii}

Add a typed interaction record to \entity{DecisionPoint}.

\begin{lstlisting}
pub struct MedicationDecisionContext {
    pub interactions_considered: Vec<MedicationInteractionRecord>,
}

pub struct MedicationInteractionRecord {
    pub interaction_type:  MedicationInteractionType,
    pub involved_fact_ids: Vec<FactId>,
    pub clinical_note:     Option<String>,
}

pub enum MedicationInteractionType {
    SideEffectManaged,
    ContraindicationOverridden,
    DoseAdjustedDueTo,
    CombinationIntentional,
    Replaced,
}
\end{lstlisting}

\begin{prosbox}
\textbf{Pros.}
\begin{itemize}[noitemsep]
  \item Keeps interaction reasoning in the clinical interpretation layer
    where it originates
  \item No new entity types
  \item Authored by a clinician confirming a narrative draft
\end{itemize}
\end{prosbox}

\begin{consbox}
\textbf{Cons.}
\begin{itemize}[noitemsep]
  \item Only captures explicitly documented interactions
  \item Nested inside \entity{Narrative} $\to$ \entity{DecisionPoint};
    cross-patient queries require traversing the narrative layer
  \item Per-episode scope; cross-episode interactions are hard to
    represent
\end{itemize}
\end{consbox}

\subsection{Comparing the Interaction Options}

The options serve different purposes and are not mutually exclusive.

Option~I makes interactions structurally queryable: ``show me all
patients on a counter-medication for metformin'' is a direct lookup.

Option~II is the lightest change and provides implicit interaction
signals via episode co-membership. It works for same-episode
interactions but breaks down for cross-episode interactions and cannot
express dose-modifying or contraindication override scenarios.

Option~III keeps interaction reasoning in the narrative layer, authored
by clinicians. It is not queryable across patients without traversing
narratives.

Options~I and~III address different concerns~--- structured
queryability vs.\ authored clinical reasoning~--- and can coexist in
the schema without conflict.

\section{The Central Question: Should Prescription Separate?}
\label{sec:separation}

The lifecycle and interaction questions are connected. The choice of
lifecycle model affects what medication interaction modeling is
possible.

Under Options~A or~B, a \entity{MedicationRelation} would link
\entity{FactId} values representing prescription \emph{events}. The
clinically meaningful relationship is between the metformin
\emph{course}~--- which may have had multiple dosage modifications
over years~--- and the semaglutide course, not between any specific
prescription event for either.

Under Option~D, \entity{MedicationRelation} links
\entity{MedicationOrder} records directly, which maps naturally to
how clinicians think about drug interactions. The prescription events
become a timeline on the \entity{MedicationOrder}, analogous to how
\entity{EpisodeMembership} events are a timeline on
\entity{ProblemEpisode}.

\begin{center}
\renewcommand{\arraystretch}{1.4}
\begin{tabular}{p{4.0cm} p{4.0cm} p{4.0cm}}
\toprule
\textbf{Concern} & \textbf{Options A/B (within Fact)} &
  \textbf{Option D (separate entity)} \\
\midrule
Schema complexity      & Low                    & High \\
Semantic correctness   & Good                   & Best \\
Temporal model         & Point per event        & \field{started\_at} / \field{stopped\_at} \\
Open-ended courses     & \field{intended\_duration} on Start & \field{stopped\_at: None} \\
Interaction modeling   & Awkward (event to event) & Natural (order to order) \\
Current med query      & Requires event scan    & Direct status lookup \\
FHIR alignment         & Partial                & Strong \\
Migration path         & Well-defined           & N/A \\
\bottomrule
\end{tabular}
\end{center}

One practical consideration: if Option~A is adopted first, migration
to Option~D is well-defined. \entity{MedicationOrder} is essentially
a materialized view of the Option~A event stream, so no clinical data
is lost in the transition. The \field{medication\_order\_id} introduced
in Option~A becomes the primary key of the \entity{MedicationOrder}
entity under Option~D.


\end{document}